% MAS1004 과제 1: 짧은 보고서, 한국어 양식.
%
% 이 과제를 하며 배운 것을 자유롭게 작성하면 됩니다.
%   - 어떤 주제를 골랐고, 왜 골랐는지
%   - 과제가 어떻게 진행되었는지
%   - 무엇을 배웠는지
%
%  이 리포트는 본문을 한쪽으로 작성하며, LLM을 쓰지 않고 직접 작성합니다.
% 본문 뒤에 그림을 세 쪽까지 더 붙일 수 있습니다.
% XeLaTeX로 컴파일하면 됩니다.  (Overleaf: Menu, Compiler, XeLaTeX).

\documentclass[11pt]{article}
\usepackage[onepage,ko]{mas1004}

\reporttitle{과제 1: Short Report}
\studentname{김민주}
\studentid{20261174}

\begin{document}
\makeheader
 저는 이 과제에서 치이카와를 주제로 골랐습니다. 
long report에서도 언급했듯이, 제가 이 주제를 선택한 이유는 '캐릭터들 서로가 유사하게 생겼음에도 잘 분류할 수 있을까?'라는 궁금증 때문이었습니다. 
물론 처음 동기는 최근 개봉한 극장판 때문이지만, 모델이 2D 캐릭터라도 어떤 지점에 집중하는지, 또한 2D 캐릭터의 굿즈도 잘 구별할 수 있는지 궁금증이 이어져 이 주제를 선정하게 되었습니다.

 과제를 진행하면서 주제에 중요한 문제가 있다는 것을 알게 되었습니다. 
모델이 검색해 모은 자료는 다수 잘못된 자료였습니다. 
찾은 자료의 문제점을 세가지로 정리하면, 첫번째는 주인공 이름과 만화의 제목이 동일하기 때문에, 검색한 자료가 주인공 외에도 다른 캐릭터를 보여주고 있었습니다. 
둘째는 다른 사람 또는 다른 만화의 캐릭터가 포함되는 것, 마지막으로 둘 이상의 캐릭터가 한 이미지에 존재하는 것이었습니다. 
이 3가지 조건으로 적절한 사진을 선별하였고, '치이카와'가 16장, '하치와레'가 112장, '우사기'가 51장으로 첫 훈련을 진행하였습니다. 그 결과 모델은 몇 개의 '치이카와' 사진을 '우사기'로 판단하였습니다. 
모델이 착각한 사진은 대부분 뒤 돌아보는 뒷모습, 치이카와의 굿즈 사진 등이었습니다. 저는 하치와레의 구별 정확도는 높은 반면 치이카와와 우사기를 구분하는 정확도는 낮은 모습에 주목했습니다. 
그리고 이 문제가 발생한 이유를 불충분한 자료 때문이라고 판단하였고, 저는 모델이 훈련하는 자료인 clean파일에 이전의 이미지와 겹치지 않는 검색 사진을 추가하였습니다. 
하지만 저는 자료를 찾는 데 너무 지쳐 두번째 훈련을 '치이카와' 38장, '하치와레' 112장, '우사기' 58장의 자료로 진행하였습니다. 
이 때 clean의 자료는 다양한 구도의 2D 캐릭터 이미지가 많았기 때문에 뒷모습의 구별력이 상승하였지만, 여전히 1차 훈련 때와 같은 굿즈 사진을 맞게 구별하지 못했습니다. 
저는 clean에게 자문을 구했고, '각 카테고리당 최소 100장을 준비하면 정확도가 올라갈 것이다'라는 말을 통해 저는 '치이카와' 104장, '하치와레' 112장, '우사기' 99장의 자료로 세번째 훈련을 진행하였습니다. 
저번 결과를 통해 추가된 사진의 일부는 굿즈 사진을 추가하여 굿즈 분별력을 높이도록 의도하였습니다. 

 100장 가량의 이미지 자료는 모델의 다양한 구도와 굿즈의 판별력을 높였습니다. 
하지만 웹사이트의 Draw it 창에서는 1차때와는 다른 결과가 발생하였습니다. 1차에서는 우사기 얼굴을 그렸을 때, 옳게 판단하였습니다. 
하지만 최종 결과인 3차에서는 Draw it에서 우사기라는 결과를 도출시키는 것이 까다로웠습니다. 
얼굴 그림은 무조건 하치와레라고 판단하고, 가끔 전신 실루엣을 보여주면 우사기라고 판단하였습니다. 
저는 이 지점에서 모델이 이미지의 어떤 부분에 집중하여 판단하는지 더 헷갈리게 되었습니다. 
왜냐하면 실제 카메라를 켜서 치이카와, 하치와레, 우사기와 관련 없는 물체를 보여주었을 때 대부분 하치와레로 인식하고, 노란색을 보이면 우사기로 판단하였습니다.
따라서 저는 모델이 색깔로 차이를 판단하는 것으로 깨닫고 각 카테고리당 자료 10개씩 흑백 사진으로 바꾸어 정확도를 확인하였습니다. 
test를 진행한 결과 30장의 사진 중 3개를 모두 치이카와로 판단하였습니다. 정확도가 매우 높을 뿐만 아니라 예상된 움직임(하치와레로 움직일 것)도 달랐습니다.
혹시 Draw it의 특징인 굵은 선 때문에 캐릭터 외형이 망가져 잘 판별하지 못한 것이 아닌지 의문이 들기 시작했고, 
앞과 같이 각 카테고리당 자료 10개씩 선의 굵기를 약 1.6배 정도(좌우 4픽셀씩 증가) 증가시켜 정확도를 확인하였습니다. 
하지만 모델은 30개 중 29개를 맞추는 놀라운 정확도를 보여주었습니다.
저는 그래서 cline에게 'Draw it에서 우사기가 하치와레로 판별된다'라는 간단한 말 대신 구체적으로 사진을 첨부하여 상황을 설명하였습니다. 
총 두장이었는데, 한장은 우사기의 얼굴만 크롭하여 그린 사진, 나머지는 우사기의 전체 실루엣만 그린 그림이었습니다. 
Draw it에서는 제가 우사기를 의도하고 그린 얼굴 그림은 항상 하치와레로 판별나고, 실루엣은 운이 좋으면 우사기, 나머지는 다 하치와레로 판단되었습니다. 
cline은 저에게 이전과는 다른 사진을 제가 선별하여 알려주면 이를 실루엣 처리를 하고 test할 것을 제안하였습니다. 이 테스트 결과 모델은 32개 중(제가 실수로 2개 더 넣었나봅니다) 17개만 옳게 판단하였습니다. 
구체적으로 말하자면 치이카와는 11장 중 1장, 하치와레는 10장 중 10장, 우사기는 11장 중 6장을 정확하게 판단하였습니다. 

 저는 이 세가지 실험을 통해 추측성 결론을 내렸습니다. 우선 색은 판단 척도에 쓰이나 결정적인 요소는 아니었습니다. 
또한 얼굴의 세부적인 가는 선을 굵게 뭉개버려도 정확도가 매우 높았습니다. 
하지만 실루엣, 그들의 귀 모양만 남은 모양 단순화는 모델의 정확도를 절반으로 낮춰버렸습니다. 
실루엣 실험에서 15장은 모두 하치와레로 이동하였습니다. 

\end{document}
